\section{Threat Model and Security Objectives}
\label{sec:threat}

\subsection{Assets}

Sirus is intended to protect four classes of assets:

\begin{itemize}
  \item \textbf{Proprietary integer logic}: license predicates, protocol
        constants, checksums, and business rules.
  \item \textbf{Embedded secrets}: salts, derivation keys, and constants that
        participate in a larger key schedule.
  \item \textbf{Algorithm structure}: the control flow and arithmetic shape that
        reveal how a decision is made.
  \item \textbf{Build geometry}: the per-build ISA assignment, register maps,
        decode-table key, and payload keys that make one protected image
        different from another.
\end{itemize}

The protected function is assumed to be a leaf or a small cluster of directly
called functions. Its result is consumed by the host as data, not as a single
boolean that an attacker can patch at the call site without consequence.

\subsection{Adversary capabilities}

We assume an adversary who can:

\begin{itemize}
  \item read the complete binary, symbol table, section headers, and metadata;
  \item run disassemblers and decompilers, including Ghidra, IDA, Binary Ninja,
        radare2, and custom parsers;
  \item run MBA simplifiers such as GAMBA, SiMBA, and MBA-Blast;
  \item run symbolic execution and SMT solvers such as angr, KLEE, Triton, and
        Z3;
  \item trace execution with Intel PIN, Frida, DynamoRIO, QEMU, gdb, or a
        hypervisor;
  \item dump process memory at any time, including during an executing
        instruction;
  \item patch file or memory bytes, including the generated entry thunk;
  \item compare two protected builds of the same source;
  \item fuzz the VM with malformed, truncated, or mutated payloads;
  \item use hardware breakpoints and debug registers.
\end{itemize}

The adversary is \emph{not} assumed to have unlimited CPU time in the intended
cost model. Sirus is an asymmetry mechanism: the defender pays a bounded
compile-time and runtime cost, and the attacker is forced through a larger
search, decoding, and analysis space. An adversary with unlimited time and full
instrumentation can eventually recover behavior; no claim in this paper
contradicts that.

\subsection{Security objectives}

Table~\ref{tab:objectives} states the objectives and the controls that address
them. The evidence column identifies the class of evidence, not a proof of
absolute security.

\begin{table}[htbp]
\centering
\small
\caption{Security objectives and current controls.}
\label{tab:objectives}
\begin{tabularx}{\textwidth}{@{}>{\raggedright\arraybackslash}p{0.06\textwidth}>{\raggedright\arraybackslash}p{0.25\textwidth}>{\raggedright\arraybackslash}p{0.35\textwidth}>{\raggedright\arraybackslash}X@{}}
\toprule
ID & Objective & Control & Evidence \\
\midrule
O1 & Remove readable original logic & Original body scrubbed; bytecode flattened, shuffled, and encrypted & Patch/scrub checks; Ghidra and radare2 output \\
O2 & Obscure operations & Bounded linear and non-linear MBA; constant blinding; randomized ISA & MBA equivalence tests; external simplifier study \\
O3 & Make builds independent & New seed changes opcodes, operand order, register maps, state ids, block order, keys, MBA shapes, decoys, and layout & Differential payload and segment equality \\
O4 & Detect payload tampering & SipHash-derived integrity value mixed into the effective payload key & Tamper experiments produce deterministic wrong output \\
O5 & Keep the host safe & Safe Rust, checked decode, no heap, fuel bound, private memory, fail-closed status & Kani, Miri, libFuzzer, malformed-payload suite, valgrind \\
O6 & Resist forced-return bypass & Protected results are data consumed downstream; wrong values change observable output & Forced-return experiment \\
O7 & Avoid leaked VM strings & Generated engine uses no formatting, panic strings, or libc symbols & Static string scan of the injected segment \\
\bottomrule
\end{tabularx}
\end{table}

\subsection{Non-objectives}

The following are explicitly outside the threat model:

\begin{itemize}
  \item protection against an adversary with unlimited local time and full
        hardware instrumentation;
  \item kernel-level anti-debugging or DBI detection in the safe tier;
  \item hiding the existence of a VM or of the protected entry points;
  \item preventing a caller-side patch that ignores the protected result;
  \item a formal cryptographic proof for the custom payload construction;
  \item whole-program virtualization;
  \item protecting functions that require external calls, host memory, floating
        point, or unsupported instructions;
  \item trusted production signing or an independent audit.
\end{itemize}

\subsection{Failure policy}

The runtime distinguishes five outcomes: \code{Returned}, \code{BadPayload},
\code{FuelExhausted}, \code{Trap}, and \code{LoopBound}. A malformed payload,
an invalid register, an out-of-range memory access, a missing manifest entry, an
invalid opcode, or an exhausted fuel budget never unwinds and never accesses host
memory. It returns value zero with the corresponding status. The build-time
pipeline similarly fails with a typed error and does not write a partial output
file.

Returning zero is not itself a security guarantee. The intended deployment
consumes the protected result as data in a downstream computation, so an
incorrect result causes an observable failure rather than a silent success.

\subsection{Cost model and evidence discipline}

We evaluate the system along three axes: correctness, safety, and recovery cost.
Correctness is measured by differential execution on real binaries. Safety is
measured by the Rust type system, bounded proofs, dynamic analysis, and fuzzing.
Recovery cost is measured by differential builds, external static and symbolic
tools, MBA simplifiers, and timing measurements. Each result is reported with
its scope. A Kani harness that proves a bounded memory helper is not presented as
a proof of the dispatcher; a Ghidra timeout is not presented as a proof that
decompilation is impossible; a dudect run below the threshold is not presented as
a formal constant-time proof.

The following terminology is used throughout:

\begin{definition}[Measured equality]
For two protected images $B'_0$ and $B'_1$ of the same input under different
seeds, the \emph{equal-byte ratio} of a section is the number of byte positions
with identical values at the same relative offset, divided by the section length.
It is a diversity metric, not a semantic security proof.
\end{definition}

\begin{definition}[Bounded symbolic resistance]
A symbolic probe is \emph{bounded-resistant} for a time budget $T$ if it does not
produce a concrete preimage for the requested target within $T$. A timeout is
evidence of cost, not a proof of unsatisfiability.
\end{definition}

\begin{definition}[Fixed-pass oblivious execution]
An execution is \emph{fixed-pass oblivious} with respect to a set of secret
inputs if, for a fixed program and pass bound $P$, the sequence of decoded basic
blocks and the sequence of private-memory addresses are independent of those
inputs, provided the active frontier converges within $P$. Register values and
branch conditions may still depend on the inputs.
\end{definition}
